\section{Introduction}\label{sec:introduction}
A subgraph of an edge-colored graph is \emph{rainbow} if its edges have pairwise distinct colors. Burr, Erd\H{o}s, Graham, and S\'os~\cite{BEGS1989} asked how few colors are needed to make every copy of a fixed graph $H$ rainbow. For a finite simple graph $G$, let $r_H(G)$ denote this minimum and, following Buci\'c, Chen, and Ma~\cite{BCM2026}, define
\begin{equation}\label{eq:fdefinition}
 f(n,e,H)=\min\set{r_H(G):\abs{V(G)}=n,\ e(G)\ge e},
\end{equation}
whenever the family of graphs is nonempty. Copies of $H$ are subgraphs, not necessarily induced; thus a copy of a cycle has distinct vertices but may have chords in $G$. Colorings need not be proper, and only colors that actually occur on an edge are counted. In~\cite{BEGS1989} the same quantity is denoted $\chi_S(n,e,H)$ and the host is required to have exactly $e$ edges. The two definitions agree, because deleting edges preserves the rainbow property and cannot increase the number of colors.

For odd cycles the natural threshold is $e=\lfloor n^2/4\rfloor+1$, one more than the Tur\'an number of $C_{2k+1}$ for large $n$, and there the answer depends strongly on the length of the cycle. For triangles the value is $3$, and for $C_5$ it is $\lfloor n/2\rfloor+3$ for large $n$, a result of Erd\H{o}s and Simonovits (see~\cite{BEGS1989,BCM2026}). For longer cycles, Burr, Erd\H{o}s, Graham, and S\'os proved that $f(n,\lfloor n^2/4\rfloor+1,C_{2k+1})$ is at least a positive multiple of $n^2$ for every fixed $k\ge3$~\cite[Theorem~5.1]{BEGS1989}. They conjectured that the correct constant is $1/8$, the value given by two disjoint cliques of nearly equal size. This conjecture is Erd\H{o}s Problem~\#809. Buci\'c, Chen, and Ma~\cite{BCM2026} proved it for every $k\ge4$. We settle the remaining case $k=3$.

\begin{theorem}\label{thm:main}
As $n\to\infty$ through all positive integers,
\begin{equation}\label{eq:main}
 f\left(n,\left\lfloor n^2/4\right\rfloor+1,C_7\right)
 =\left(\frac18+o(1)\right)n^2.
\end{equation}
\end{theorem}

The upper bound comes from the two-clique construction. The new content is the lower bound: for every $\varepsilon>0$ and all sufficiently large $n$, every $n$-vertex graph with at least $\lfloor n^2/4\rfloor+1$ edges needs at least $(1/8-\varepsilon)n^2$ colors in any edge-coloring in which every $C_7$ is rainbow.

\begin{corollary}\label{cor:allodd}
For every fixed integer $k\ge3$,
\begin{equation}\label{eq:allodd}
 f\left(n,\left\lfloor n^2/4\right\rfloor+1,C_{2k+1}\right)
 =\left(\frac18+o(1)\right)n^2.
\end{equation}
The case $k=3$ is Theorem~\ref{thm:main}; the cases $k\ge4$ are due to Buci\'c, Chen, and Ma.
\end{corollary}
\begin{proof}
For $k\ge4$, apply~\cite[Theorem~1.2]{BCM2026} at $e=\lfloor n^2/4\rfloor+1$; its square-root term is $O(n)$. The remaining case is Theorem~\ref{thm:main}.
\end{proof}

\subsection{Outline of the proof}
The \emph{conflict graph} of $G$ has vertex set $E(G)$, two edges being adjacent when they lie on a common simple $C_7$; its chromatic number is $r_{C_7}(G)$. Rather than work with this graph directly, we pass to a vertex-weighted graph and replace common seven-cycles by a condition on walks: two edges are \emph{compatible} if their endpoints are not joined by complementary walks of lengths two and three. After a small number of edges have been deleted, a path-lifting lemma turns such walks into simple paths of the original graph, and the original color classes then satisfy the compatibility relation.

An edge is \emph{triangular} if it lies in a triangle. Once the weighted minimum degree exceeds $1/3$, a set of pairwise compatible edges is either entirely triangular or entirely nontriangular, and in the first case it has at most two edges. The saving available in the triangular part is therefore the value $\nu$ of a capacitated fractional matching. In the nontriangular part, the geometry of neighborhoods gives a family of lower bounds $q(w)$ on the palette cost. The heart of the argument is the inequality
\[
 2\sum_wx_w\tau(w)q(w)
 +\left(\frac12-m-2f-2\nu\right)\sum_wx_w\tau(w)\ge0,
\]
where $m$ is the total edge mass, $f$ is the nontriangular edge mass, and $\tau(w)$ is the edge mass inside $N(w)$. We prove it with an exact rational certificate on five sampled vertices. The certificate uses two constraints supplied by the triangular matching: a bound on endpoint incidences in each neighborhood, and a bound on the union of the neighborhoods of each matched edge.

The resulting palette bound has a stable form that tolerates $m$ slightly below $1/4$. Some such stability is necessary. Path lifting relies on regularity and triangle removal and therefore deletes edges, while the surplus of the original graph over $n^2/4$ may be a single edge. When every bipartite cut misses many edges, a lower bound on the triangle mass pays for the deleted edges. When some cut contains almost all edges, we find a large clique in the conflict graph directly. Finally, a vertex-deletion argument removes the minimum-degree assumption while keeping the number of edges strictly above $h^2/4$, where $h$ is the current order.

Sections~\ref{sec:palettes}--\ref{sec:finite} prove the finite palette inequality. Sections~\ref{sec:cleanup}--\ref{sec:completion} transfer it to arbitrary graphs and complete the proof of Theorem~\ref{thm:main}. Section~\ref{sec:verification} describes the Lean formalization of Corollary~\ref{cor:allodd}, which combines the new $C_7$ proof with a formalization of the Buci\'c--Chen--Ma argument for $k\ge4$.
